
\chapter{Introduction}

\section{Motivation}

The emergence of blockchain technology has catalyzed a fundamental transformation in financial infrastructure, giving rise to a new class of decentralized financial systems that operate without traditional intermediaries. As of early 2026, global DeFi Total Value Locked (TVL) stands at approximately \$125 billion — a figure that, despite a seasonal correction from the mid-2025 peak of \$2.6 billion on Sui alone, reflects the enduring adoption and long-term growth trajectory of on-chain finance. Within this ecosystem, decentralized exchanges (DEXs) have emerged as the primary venue for permissionless asset trading, sustaining an average global daily volume of approximately \$19 billion throughout 2025 and into 2026, with Sui-native DEXs contributing approximately \$250 million daily \cite{SuiQ220252025}.

As highlighted in the OKX State of DEXs 2025 Report, the central promise of decentralized trading is the delivery of low slippage and best-price execution. Since the inception of DEXs in 2018, successive architectural generations have pursued this goal through distinct mechanisms. Automated Market Makers (AMMs) democratized liquidity provision but introduced price impact proportional to trade size. Order book DEXs improved price precision but required active market participation and remained susceptible to liquidity fragmentation. DEX Aggregators improved execution by routing across multiple venues, yet remained fundamentally constrained by the depth of underlying pools. Hybrid models combine these approaches at the cost of architectural complexity. Each generation represents an incremental improvement; none fully resolves the slippage problem under fragmented liquidity conditions.

It is in this context that Maximal Extractable Value (MEV) emerges as a critical phenomenon. MEV refers to the profit that block producers and sophisticated market participants can extract by strategically reordering, including, or excluding transactions within a block prior to finalization \cite{WhatMaximalExtractable}. This is made possible by the public visibility of pending transactions in the mempool: before a trade is confirmed, its parameters are observable by any network participant. Searchers — automated agents monitoring the mempool — exploit this transparency to execute three primary forms of MEV extraction. In \textbf{front-running}, a searcher identifies a profitable pending trade and submits an identical transaction with a higher priority fee, ensuring execution ahead of the original user and capturing the resulting price movement. In a \textbf{sandwich attack}, the attacker places a buy order immediately before the target transaction and a sell order immediately after, artificially inflating the execution price for the victim while profiting from the induced price delta. In \textbf{back-running}, the attacker inserts a transaction directly following a price-moving event to capture the resulting arbitrage opportunity before the market equilibrates \cite{WhatMaximalExtractable}.

The implications of MEV extend beyond individual transaction losses. At scale, MEV introduces systemic inefficiencies: it increases effective transaction costs for ordinary users, distorts price formation on DEXs, and creates an asymmetric information advantage for sophisticated actors with access to specialized infrastructure. These effects collectively undermine the fairness and accessibility that decentralized finance purports to guarantee \cite{WhatMaximalExtractable}.

In response to MEV's growing impact, multiple mitigation approaches have emerged. Flashbots introduced private transaction relaying to reduce mempool exposure; ZeroMEV developed analytics tooling to quantify and attribute extracted value; and CoW Protocol introduced a structurally distinct execution model. Proposed in 2021, CoW Protocol's \textbf{Coincidence of Wants} mechanism batches user trade intents and searches for peer-to-peer matches — settling directly between counterparties before routing residual volume to external liquidity sources. When a user intending to swap token A for token B is matched with a counterparty intending the inverse trade, both orders are settled at the reference market price with zero AMM fee and zero price impact. This batch auction model simultaneously eliminates MEV exposure, reduces slippage, and recovers surplus for users rather than extractors. According to on-chain data from Dune Analytics, even on Ethereum the average CoW batch contains approximately 1.5 intents, illustrating that peer-to-peer matching remains probabilistic and volume-dependent \cite{BatchHistoryDune2023}.

On Sui, MEV is addressed through a different set of architectural mechanisms — including priority gas auctions (PGAs) and consensus amplification via SIP-45 — with the explicit design principle that user transaction protection takes priority over extracted value \cite{MEVSuiCurrent2025}. As a result, the MEV problem that primarily motivates CoW on Ethereum is partially mitigated by Sui's protocol design. Nevertheless, the CoW model retains a powerful and distinct advantage: the ability to achieve \textbf{zero slippage} when sufficient matching volume exists. Furthermore, a CoW-style protocol on Sui does not compete with or displace existing DEXs. Rather, it functions as a \textbf{meta-DEX aggregator} that discovers and exploits better pricing across Cetus, Aftermath, DeepBook, and other Sui protocols — reducing liquidity fragmentation rather than adding to it. As confirmed by the Sui Q2 2025 DeFi Roundup \cite{SuiQ220252025}, AMMs and aggregators currently dominate Sui DeFi volume, and no major CoW-style batch auction protocol exists on the network, representing a clear gap in the ecosystem.

Understanding the behavior, scale, and distribution of MEV — as well as the effectiveness of countermeasures such as CoW-style batch auctions — has become essential for designing fair and efficient decentralized trading systems. However, existing implementations of CoW Protocol on Ethereum rely on off-chain infrastructure for solver coordination, intent relay, and price verification, reintroducing trusted intermediaries whose failure modes have been demonstrated in practice: a solver network exploit resulted in direct fund losses, and a separate incident in which the off-chain permission model caused a user to lose nearly \$50 million in a misrouted trade — both illustrating the systemic risk of off-chain trust assumptions \cite{sinclairCryptoTraderLoses2026}.

This thesis addresses the gap between the theoretical promise of the CoW model and its current implementation constraints. Leveraging Sui's object-centric architecture — which provides native primitives for shared object ownership, programmable transaction blocks, and sub-second finality — this work designs and implements a trustless CoW batch auction protocol in which the entire lifecycle, from intent locking through solver competition to atomic on-chain settlement, is enforced by smart contract logic with no off-chain coordinator. The result is a meta-DEX aggregator that complements Sui's existing liquidity ecosystem, reduces fragmentation, and provides users with best-execution guarantees backed by on-chain price verification.

\section{Problem Statement}

Although the Coincidence of Wants model has demonstrated its effectiveness in optimizing trade execution and reducing slippage on Ethereum, its application within the Sui ecosystem remains largely unexplored. This thesis identifies four core problems that motivate the proposed design.

\textbf{P1 --- Absence of a CoW-based Protocol on Sui}

Despite the significant growth of Sui's DeFi ecosystem --- with total DEX volume exceeding \$33 billion in Q2 2025 alone \cite{SuiQ220252025} --- no major protocol operating on the CoW batch auction model currently exists on the network. Existing protocols, including AMMs such as Cetus, Bluefin, and Aftermath, and aggregators such as 7k and FlowX, operate independently and lack a coordination layer capable of identifying and exploiting coincidence-of-wants opportunities across concurrent trade intents. As a consequence, liquidity remains fragmented across isolated pools, and users continue to incur unnecessary slippage even in market conditions where sufficient opposing intent volume exists. The absence of a CoW-style meta-DEX aggregator means the Sui ecosystem foregoes the possibility of zero-slippage execution in eligible trade conditions, while leaving unexploited the natural alignment between Sui's object-centric architecture and the intent-based trading model.

\textbf{P2 --- Off-chain Trust Assumptions and Permission Risk in Existing CoW Implementations}

The architecture of CoW Protocol on Ethereum relies substantially on off-chain infrastructure for core protocol functions, including intent relaying, solver coordination, and price verification. This design reintroduces dependence on trusted intermediaries --- a dependency that decentralized finance is fundamentally intended to eliminate. The practical consequences of this trust assumption have been demonstrated through two documented incidents. First, an exploit targeting CoW Protocol's solver network resulted in direct user fund losses \cite{cowswapsolverexploit}. Second, a separate incident in which the off-chain permission and approval model caused a user to lose approximately \$50 million in a misrouted trade --- receiving only a \$600,000 fee refund as compensation \cite{sinclairCryptoTraderLoses2026}. Both cases illustrate that when critical execution decisions are processed off-chain, users cannot rely on smart contract guarantees to protect their assets. A genuinely trustless system requires that the entire transaction lifecycle --- intent locking, solver competition, price verification, and settlement --- be enforced and verifiable on-chain.

\textbf{P3 --- High Operational Cost and Solver Capital Inefficiency}

On Ethereum-based deployments, high gas costs impose two structural constraints on the efficiency of the CoW model. First, per-transaction costs force solvers to accumulate a sufficient number of intents before a batch becomes economically viable, resulting in extended user wait times. In practice, CoW Protocol on Ethereum processes batches approximately every ten minutes \cite{BatchHistoryDune2023}, during which users have no execution guarantee. Second, solvers must maintain substantial capital reserves to fund settlement for complex batches, creating high barriers to entry and limiting competition within the solver network. On Sui, significantly lower gas fees and sub-second transaction finality open the possibility of high-frequency batch processing --- enabling what may be characterized as micro-batching. Furthermore, Sui's Programmable Transaction Block (PTB) mechanism allows solvers to source on-chain liquidity from DeepBook within the same atomic transaction as settlement, reducing capital requirements to near zero and lowering the barrier for solver participation.

\textbf{P4 --- Liquidity Fragmentation and the Absence of an Optimal Routing Layer}

Even in the absence of direct peer-to-peer matching opportunities, users on Sui face the challenge of identifying optimal execution paths across fragmented liquidity venues. Existing aggregators address this partially by routing individual orders across multiple pools, but they operate reactively --- processing each order in isolation as it arrives, without the ability to combine concurrent intents and exploit potential CoW matches within a shared time window. As a result, a meaningful portion of trades continues to incur unnecessary slippage by being routed through AMM pools when a directly opposing intent, submitted within the same batch window, could have fulfilled the order at a better price. A meta-aggregator built on the batch auction model can address both dimensions simultaneously: maximizing the probability of peer-to-peer CoW matching and guaranteeing optimal AMM routing for residual unmatched volume.

% \section{Research Objectives}

% \section{Contributions}